\subsection{同态模的算子环的扩张}

设A为交换环，B为环， \(\rho: \mathbf{A} \to \mathbf{B}\) 为环同态，E、F为两个A-模；由于B（借助 \(\rho\)）是(A,A)-双模，且F可视为(A,A)-双模，因此在Z-模 B \(\otimes_{\mathbf{A}}\) F 上有两种A-模结构，在这些结构下分别有 \(a(b \otimes y) = (\rho(a)b) \otimes y\) 和 \(a(b \otimes y) = b \otimes (ay)\)（对 \(a \in \mathbf{A}\)、\(b \in \mathbf{B}\)、\(y \in \mathbf{F}\)）。我们把这样定义的两个A-模记为G'和G''；此外，G'恰好就是A-模 \(\rho_*(\rho^*(\mathbf{F}))\)。

这样一来，在§4第2号公式(7)的典范同态的定义中，我们把B换成A，把B-模F换成借助 \(\rho\) 视为A-模的环B，并把G换成视为(A, A)-双模的F；由于A交换，所得的典范Z-同态可写作

\[
\mathbf {B} \otimes_ {\mathbf {A}} \operatorname{Hom} _ {\mathbf {A}} (\mathbf {E}, \mathbf {F}) \rightarrow \operatorname{Hom} _ {\mathbf {A}} (\mathbf {E}, \mathbf {G} ^ {\prime \prime}).\tag{15}
\]

另一方面（第1号，公式(2)），存在一个典范Z-同构

\[
\mathrm{Hom} _ {\mathbf {A}} (\mathbf {E}, \mathbf {G} ^ {\prime}) = \mathrm{Hom} _ {\mathbf {A}} (\mathbf {E}, \rho_ {*} (\rho^ {*} (\mathbf {F}))) \rightarrow \mathrm{Hom} _ {\mathbf {B}} (\rho^ {*} (\mathbf {E}), \rho^ {*} (\mathbf {F})). \tag {16}
\]

现在假设 \(\rho(A)\) 包含于 \(B\) 的中心，这时 \(\rho\) 也称为中心同态 \(*\)（或称 \(\rho\) 在 \(B\) 上定义了一个A-代数结构，参见III，§1，第3号）。则 \(G'\) 与 \(G''\) 的A-模结构相同，并且把同态（16）与（15）复合，我们从而得到一个典范Z-同态

\[
\omega : \mathrm{B} \otimes_ {\mathrm{A}} \operatorname{Hom} _ {\mathrm{A}} (\mathrm{E}, \mathrm{F}) \rightarrow \operatorname{Hom} _ {\mathrm{B}} \left(\mathrm{E} _ {(\mathrm{B})}, \mathrm{F} _ {(\mathrm{B})}\right)\tag{17}
\]

其特征是，对于所有 \(u \in \operatorname{Hom}_{\mathbf{A}}(\mathbf{E}, \mathbf{F})\) 以及所有 \(b \in B\)

\[
\omega (b \otimes u) = r _ {b} \otimes u,\tag{18}
\]

其中 \(r_b\) 表示B中右乘 \(b\) 的映射。

此外，\(\omega\) 是中心同态这一假设蕴含 \((bb')\rho(a) = b\rho(a)b'\)（对B中的 \(b, b'\) 及 \(a \in A\)）；换言之，\(B_d\) 的右B-模结构与其A-模结构相容；从而它在 B \(\otimes_A\) \(\mathrm{Hom}_{\mathbf{A}}(\mathbf{E}, \mathbf{F})\) 上定义了一个右B-模结构（§3，第4号），也在 \(\mathrm{F}_{(\mathbf{B})} = \mathbf{B} \otimes_{\mathbf{A}} \mathbf{F}\) 上定义一个；最后，由于 \(\mathrm{F}_{(\mathbf{B})}\) 上的左、右B-模结构相容，在 \(\mathrm{Hom}_{\mathbf{B}}(\mathrm{E}_{(\mathbf{B})}, \mathrm{F}_{(\mathbf{B})})\) 上也得到一个右B-模结构（§1，第14号）。然后立即验证（17）对这些结构是右B-模同态。

命题7. 设 A 是一个交换环，B 是一个环， \(\rho: \mathbf{A} \to \mathbf{B}\) 一个中心同态，且 E、F 是两个 A-模。

\begin{enumerate}
\def\labelenumi{(\roman{enumi})}
\item
  若 B 是投射（相应地，有限生成投射）A-模，则同态 (17) 是单射（相应地，双射）。
\item
  若 \(\mathbf{E}\) 是有限生成的投射A-模，则同态 (17) 是双射。
\end{enumerate}

由于（16）是双射，该命题由§4第2号命题2应用于典范同态（15）得出。

